\section{Introduction}
\label{sec:intro}

As large language models (LLMs) continue to advance at an exponential pace, their range of applications has expanded rapidly, accompanied by growing concerns about their safety and reliability.
To prevent models from generating harmful content, researchers have invested heavily in safety alignment through reinforcement learning from human feedback (RLHF) and safety fine-tuning~\cite{rlhf, safetyft}.
Despite these defenses, malicious users craft jailbreak attacks that bypass safety policies and steer models into producing content that violates them.

\begin{figure}[t]
      \centering
      \includegraphics[width=\textwidth]{architecture.png}
      \caption{Overview of the IAM-Teamer framework. A harmful goal sampled from
            AdvBench is sent to the Attack Model, which generates an attack prompt;
            the Target Model produces a response, and the Judge Model scores it
            along five axes (off-topic, refusal, convincing, specific, toxicity) to
            yield a score and a set of success factors. Each attempt (attack prompt,
            target response, score, strategy, success factor, and refusal pattern)
            is recorded in the MemoryDB, whose accumulated target-specific feedback
            conditions the Attack Model's next prompt.}
      \label{fig:architecture}
\end{figure}

Research on such attacks has advanced in step with these defenses.
Early jailbreaks relied on hand-crafted single prompts (e.g., token- or character-level perturbations, persona play, fictional framing, or prefix injection).
However, these static and predictable methods are straightforward to defend against.
Stronger white-box optimizers instead search for adversarial inputs using gradient signals~\cite{gcg}, but they require model internals that deployed systems do not expose.
Recent work has introduced automated frameworks, such as Crescendo~\cite{crescendo} and GOAT~\cite{goat}, to orchestrate multi-turn jailbreak conversations.
We review this progression in detail in Section~\ref{sec:related}.

Despite this progress, even the most recent agentic frameworks share two limitations.
First, the prompts they generate lack \emph{diversity}, collapsing into a narrow set of recurring patterns.
Second, they apply \emph{a single strategy} uniformly to all models, regardless of their differing characteristics and safety alignments.

To address these limitations, we propose IAM-Teamer (Inference-Adaptive Memory for Autonomous Red Teamer), a framework built from three models (Figure~\ref{fig:architecture}): an attack model that generates prompts, a target model that responds, and a judge model that scores each response.
Every attempt (its prompt, the strategy used, and the resulting score) is recorded in a memory maintained separately for each target, and the attacker conditions its next prompt on this accumulated, target-specific feedback.
By leveraging this per-target memory, IAM-Teamer sustains strategy exploration while adapting to specific model defenses, thereby achieving successful jailbreaks in substantially fewer turns.
A comprehensive comparison with representative frameworks is summarized in Table~\ref{tab:framework_comparison}.

\begin{table}[htbp]
      \centering
      \caption{Comparison with existing jailbreak frameworks.}
      \label{tab:framework_comparison}
      \footnotesize
      \renewcommand{\arraystretch}{1.3}
      \setlength{\aboverulesep}{0pt}
      \setlength{\belowrulesep}{0pt}
      \setlength{\tabcolsep}{6pt}
      \begin{tabular}{lccccc}
            \toprule
            \thead[l]{Framework}                                             & \thead{Multi-turn} & \thead{Gradient-based\\optimization} & \thead{Long-life\\Memory} & \thead{Target-specific\\strategy\\selection} & \thead{Strategy-\\selection\\Refinement} \\
            \midrule
            GCG~\cite{gcg} / AutoPrompt~\cite{autoprompt} &                    & $\checkmark$                         &                           &                                              &                                          \\
            PAIR~\cite{pair} / TAP~\cite{tap}                                & $\triangle$        &                                      &                           &                                              & $\triangle$                              \\
            Rainbow Teaming~\cite{rainbow}                &                    &                                      & $\triangle$               &                                              & $\triangle$                              \\
            AutoDAN-Turbo~\cite{autodanturbo}                                & $\triangle$        &                                      & $\checkmark$              &                                              &                                          \\
            AutoRedTeamer~\cite{autoredteamer}            & $\triangle$        &                                      & $\checkmark$              &                                              &                                          \\
            ASTRA~\cite{astra}                                               & $\checkmark$       &                                      & $\checkmark$              &                                              & $\triangle$                              \\
            Crescendo~\cite{crescendo}                    & $\checkmark$       &                                      &                           &                                              &                                          \\
            GOAT~\cite{goat}                                                 & $\checkmark$       &                                      &                           &                                              & $\checkmark$                             \\
            \midrule
            \textbf{IAM-Teamer (Ours)}                    & $\checkmark$       &                                      & $\checkmark$              & $\checkmark$                                 & $\checkmark$                             \\
            \bottomrule
      \end{tabular}

      \begin{minipage}{\linewidth}
            \footnotesize $\checkmark$: fully supported; $\triangle$: partially supported (e.g., iterative single-turn rather than a sustained multi-turn conversation); blank: not supported. \emph{Multi-turn}: drives a single evolving conversation. \emph{Long-life Memory}: persists learned strategies across goals. \emph{Target-specific strategy selection}: chooses strategies conditioned on the specific target's observed behavior. \emph{Strategy-selection Refinement}: updates the strategy choice from per-turn feedback.
      \end{minipage}
\end{table}

Our contributions are summarized as follows:
\begin{itemize}
      \item \textbf{Target-Adaptive Framework:} We propose IAM-Teamer, a jailbreak framework that utilizes a persistent, per-target memory to dynamically adapt attack prompts to specific model defenses.
      \item \textbf{Strategic Exploration \& Distillation:} We design a strategy pipeline that operationalizes this memory by integrating Thompson sampling for balanced strategy exploration and an automated judge for distilling successful attack mechanisms.
      \item \textbf{Empirical Efficiency:} Across six open-weight and proprietary targets, we show that IAM-Teamer matches or exceeds strong multi-turn baselines (Crescendo and GOAT), reaching success in substantially fewer turns on highly aligned targets, at a small early-turn cost on the least-aligned target where an overt first turn already succeeds quickly.
\end{itemize}
